\section{INTRODUCTION} \label{intro}
%The next line produces an indented paragraph to start the document
 %unit.  The LaTeX defaults start most units without indentations.
\hspace{\parindent}
This document has been prepared to serve as a template for preparation 
of a master's or Ph.D. thesis in mathematics at New Mexico State University.  
The template uses the \LaTeX\ $2_{\varepsilon}$  document preparation
system as the basic tool. Before you attempt to use the template, you
should become familiar with the \LaTeX\ $2_{\varepsilon}$ system.  
For example, you may browse the directory 
{\tt usr/local/teTeX/texmf/doc/latex} 
on the Department of Mathematical Sciences UNIX server for essential
information about \LaTeX\ $2_{\varepsilon}$.  You may also wish to
invest in a copy of Leslie Lamport's user guide and reference manual
\cite{Lamport}.

However, the \LaTeX\ $2_{\varepsilon}$ document preparation system by 
itself is not sufficient for preparing an NMSU thesis.  On one hand, 
some of the formatting requirements of the NMSU graduate school
conflict with the default formatting provided by the document class 
{\tt article} provided by the standard \LaTeX\ $2_{\varepsilon}$ package.
Some of the NMSU formatting requirements may be fulfilled by settings
made in the preamble of your \LaTeX\ document.  To handle other 
parts of the graduate school's formatting guidelines, 
we have provided a new document class {\tt nmsuth01} that
seems to solve most of these local formatting issues.  You will need to have 
the file {\tt nmsuth01.cls} in the directory with the files comprising 
your thesis in order for \LaTeX\ to find it.  Note that when the file 
{\tt nmsuth01.cls} loads, \LaTeX\ reports a warning that an unusual
document class is in use.  This is normal and should not effect the
creation of the {\tt .dvi} file.   Finally, we should warn potential
users that we have constructed the file {\tt nmsuth01.cls} to handle
issues that seem to be essential to us.  It is possible that
``extreme'' users will have formatting problems that necessitate
additional modifications to the article class file.

After you have gained some experience with \LaTeX\ and with literature 
in your speciality, it will become clear that certain demands placed
on the typesetter by higher mathematics cannot be fulfilled by the 
\LaTeX\ $2_{\varepsilon}$
package itself.  To get around these problems, 
users of \TeX\ have developed developed special packages 
to make the typesetting of mathematics easier.  
The American Mathematical Society has a package of macros to handle
such issues as formatting for statements of theorems, definitions,
examples, remarks, and so on.  Other parts of the American Mathematic
Society package deal with complicated  display alignment and line
breaking issues.  Other packages have been developed to handle
inclusion of graphics into mathematical documents.  
All of the packages used in the template are installed on the
Department of Mathematics UNIX server.  For general information about
what \TeX\ features are currently available browse the directory
{\tt usr/local/teTeX/texmf/doc} on the UNIX server.
Your thesis advisor should also be willing to draw your attention to
specific packages you might find especially useful.


This files for this template were prepared by Pat Morandi and Ross Staffeldt.
All the {\tt .tex} files for the template were prepared using the
emacs editor.  

Before you start creating files and problems for yourself, you should
acquire some basic background information about writing in the
mathematical sciences.  A basic reference is Nicholas J. Higham's book
\cite{Higham}, which includes discussion of what differentiates good
writing from poor writing, preparation of talks and posters, as well
as preparation of theses and papers.  He also has material on 
\TeX\ systems and a brief introduction to the emacs editor.  
The American Mathematical Society publication 
\cite{Swanson} has general information about mathematical typesetting
which is still useful, although the book predates the \TeX-age.
We conclude with a brief description of the contents of this template.

Chapter \ref{homology} is an extract from the 1997 Ph.D. thesis of
Eduardo Qui\~nones-Rico.  It is included to illustrate the some of the 
features made available in the packages 
{\tt amsmath} and {\tt amsthm}
which are in the macro package AMS-\LaTeX\ distributed free of charge
by the American Mathematical Society.  Among the features illustrated
in this section are the environments for theorems, proofs, lemmas,
remarks, and examples of ``cases'' constructions.   This section also 
illustrates some commutative diagrams illustrating the package {\tt xypic}.
In later versions of the template we hope to include examples from
other NMSU theses illustrating problems and solutions associated to
communicating other types of mathematics.

Chapter \ref{table} provides a sample table containing mathematical
entries.  The table is also extracted from Eduardo's thesis.  There
are no special packages called for in this section.  The main purpose
of the section is to give a reason for creating a list of tables
following the table of contents page.

Chapter \ref{graphics} provides two figures which are the realizations 
of two encapsulated PostScript files.  Creating and incorporating
graphics is rather sophisticated business.  For detailed information
on the incorporation of graphics into a \LaTeX\ $2_{\varepsilon}$ document 
go to some directory under your UNIX account and use the command 
\newline
{\tt cp /usr/local/teTeX/texmf/doc/latex/graphics/epslatex.ps
  epslatex.ps}
\newline
to transfer a copy of the document 
``Using EPS Graphics in \LaTeX\ $2_{\varepsilon}$ Documents''
to your account, where you can view it or print it for later reference.



